\documentclass[11pt]{article}
\usepackage[T1]{fontenc}
\usepackage[margin=1.1in]{geometry}
\usepackage{lmodern}
\usepackage{booktabs}
\usepackage{hyperref}
\usepackage{leanquote}
% The example file is a copy outside its Lake project, so it cannot be verified here.
\leanquotesetup{project=../example/, verify=false}

\lstdefinestyle{latexsrc}{language=[LaTeX]TeX, basicstyle=\ttfamily\small, columns=fullflexible,
  keepspaces=true, frame=single, rulecolor=\color{black!25}, xleftmargin=1em, xrightmargin=1em,
  moretexcs={leanfile,leanquote,leanfilequote,leanquotesetup}}
\lstnewenvironment{latexsrc}{\lstset{style=latexsrc}}{}
\newcommand{\meta}[1]{\ensuremath{\langle}\textit{#1}\ensuremath{\rangle}}
\newcommand{\cs}[1]{\texttt{\textbackslash #1}}
\newcommand{\opt}[1]{\texttt{#1}}

\title{The \textsf{leanquote} package\\[0.3em]
  \large Quote Lean declarations in \LaTeX{} by name}
\author{Stefan Langerman \and Erik D. Demaine}
\date{Version 0.8, 2026-10-25\\\url{https://github.com/slef/leanquote}}

\begin{document}
\sloppy
\maketitle

\begin{abstract}
\noindent
\textsf{leanquote} typesets declarations of a Lean~4 file in a \LaTeX{} document by name:
\cs{leanfile}\verb|{Abstract.lean}| followed by \cs{leanquote}\verb|{main_theorem}| prints the
declaration \texttt{main\_theorem} exactly as it appears in the Lean source, read at compile time.
The prose stays in \LaTeX, and the code is never retyped or pasted.
A misspelled name is an error, and the package warns when the Lean file has changed since Lean
last compiled it.
It needs nothing on the Lean side beyond the usual \texttt{lake build}; an optional exporter lets
Lean itself report the exact extent of each declaration.
It works with pdf\LaTeX, Lua\LaTeX, and Xe\LaTeX.
\end{abstract}

\section{Why}

Papers that use a formalization often reproduce its definitions and main theorem statements,
so that a reader can check that the formal statements say what the paper claims.
Pasting code by hand invites the copy in the paper to drift from the code that Lean checked.
With \textsf{leanquote} the paper names the declarations it wants, and \LaTeX{} copies them from
the Lean file each time the paper is compiled.
The package was written for papers that present a short, human-readable
\texttt{Abstract.lean}, containing the definitions and main statements of a formalization and
proved through a bridge to the full development, but it works with any Lean~4 file.

\section{Usage}

\begin{latexsrc}
\usepackage{leanquote}
\leanquotesetup{project=../lean/}            % root of the Lake project (optional)

\leanfile{TranslationTiling/Abstract.lean}  % current file
A prototile is a finite list of points of $\mathbb{Z}^3$:
\leanquote{Point, Prototile}
\end{latexsrc}

\noindent produces (with the file shipped in the package's \texttt{example} directory):

\leanfile{TranslationTiling/Abstract.lean}
A prototile is a finite list of points of $\mathbb{Z}^3$:
\leanquote{Point, Prototile}

\medskip
\begin{description}
\item[\cs{leanfile}\texttt{[\meta{options}]\{\meta{file}\}}]
  sets the current file for the following \cs{leanquote} commands. Options given here apply to
  all of them. The path is relative to \opt{project}, if set. The same command is also
  available as \cs{leanquotefile}. If the document already defines \cs{leanfile} (some
  blueprint macro files do), \textsf{leanquote} leaves that definition alone, warns, and
  \cs{leanquotefile} is used instead.
\item[\cs{leanquote}\texttt{[\meta{options}]\{\meta{names}\}}]
  typesets the comma-separated declarations \meta{names} of the current file, together, in one
  listing.
\item[\cs{leanfilequote}\texttt{[\meta{options}]\{\meta{file}\}\{\meta{names}\}}]
  does the same for an explicitly given file.
\item[\cs{leanquotesetup}\texttt{\{\meta{options}\}}]
  sets options for the rest of the document (or the current group).
\end{description}

\noindent Names may be short (\texttt{Tiles}) or qualified
(\texttt{TranslationTiling.Abstract.Tiles}). A qualified name matches a declaration inside the
corresponding \texttt{namespace}: either name may be a dotted suffix of the other. An unknown name
is an error; an ambiguous one is a warning, and the first match is used.

With the \opt{doc} option, the doc comment is included:

\begin{latexsrc}
\leanquote[doc]{Tiles}
\end{latexsrc}
\leanquote[doc]{Tiles}

\section{Options}

\begin{center}
\begin{tabular}{@{}lll@{}}
\toprule
Option & Default & Meaning \\
\midrule
\opt{doc} & false & also show each declaration's doc comment \\
\opt{box} & true & put the listing in a light frame (\texttt{leanquotebox}) \\
\opt{logo=}\meta{file} & none & image in the frame's top-right corner \\
\opt{style=}\meta{name} & \texttt{leanquote} & listings style used for the code \\
\opt{project=}\meta{dir}\texttt{/} & empty & root of the Lake project \\
\opt{olean=}\meta{file} & & compiled file to compare with, if not at the standard place \\
\opt{verify} & true & check that Lean compiled the code; if not, mark it \textsc{unverified} \\
\opt{root=}, \opt{source=} & & for explicit exports: where the recorded source is \\
any other key & & passed to \cs{lstinputlisting} \\
\bottomrule
\end{tabular}
\end{center}

\section{How declarations are found}

The package reads the Lean file with \LaTeX's own file input and uses the standard layout of
Lean code. A declaration starts at a line in column~0 of the form
\begin{center}
\texttt{[@[\meta{attributes}]] [\meta{modifiers}] \meta{keyword} \meta{name} \dots}
\end{center}
where \meta{keyword} is one of \texttt{def}, \texttt{theorem}, \texttt{lemma}, \texttt{abbrev},
\texttt{structure}, \texttt{inductive}, \texttt{instance}, \texttt{class}, \texttt{axiom},
\texttt{opaque}, and the modifiers are \texttt{private}, \texttt{protected},
\texttt{noncomputable}, \texttt{partial}, \texttt{unsafe}, \texttt{nonrec}.
Attribute lines just above it are included, and so is the doc comment if \opt{doc} is set.
The declaration continues while lines are blank, indented, or start with \texttt{|},
\texttt{termination\_by}, \texttt{decreasing\_by} or \texttt{deriving}; trailing blank lines are
dropped. This matches the layout used by Mathlib and by most Lean code, including
\texttt{match} alternatives and \texttt{inductive} constructors.

\section{Was the code compiled?}

\LaTeX{} does not run Lean. But Lake writes a module's compiled file
(\texttt{.lake/build/lib/lean/\dots/File.olean}) only when it builds successfully, so after a
successful \texttt{lake build} the \texttt{.olean} is newer than its \texttt{.lean} source.
With \opt{project} set to the root of the Lake project, \textsf{leanquote} compares the two
modification dates. If the source is newer, or if there is no compiled file to compare with,
the quote is marked \textsc{unverified} in the corner of its box, since a warning in the log is
easily missed:

\leanquote[verify]{Adjacent}

\noindent (Here the example file has no compiled \texttt{.olean}.) Without a box, a red
\textsc{unverified} line appears above the code. The log also explains the problem:
\begin{latexsrc}
Package leanquote Warning: ../lean/TranslationTiling/Abstract.lean is newer
(leanquote)  than its compiled file ... The code shown may not be what Lean
(leanquote)  checked (the quote is marked UNVERIFIED): run lake build.
\end{latexsrc}
For a copy of a Lean file outside its project, which cannot be checked, set \opt{verify=false}
to turn off both the mark and the warning; the old name \opt{check} is also accepted.
The mark's appearance can be changed by redefining \cs{leanquoteunverifiedbadge} (in the box)
and \cs{leanquoteunverifiedtext} (without a box).

\paragraph{The logo.} With \opt{logo=}\meta{file}, the image appears in the top-right corner of
each box, followed by \cs{leanquotelogomark}, a small ``TM'' in the blue of the Lean logo.
The Lean trademark policy (\url{https://lean-lang.org/trademark-policy}) asks that the Lean logo
always be accompanied by a trademark symbol; for other logos, redefine
\cs{leanquotelogomark} as empty. The Lean logo is not included in the package.

\section{Optional: exact boundaries from Lean}

For files laid out unusually, Lean itself can report where each declaration starts and ends.
From the root of the Lake project, run the exporter shipped with the package:
\begin{latexsrc}
path/to/leanquote/leanquote-export TranslationTiling.Abstract [More.Modules ...]
\end{latexsrc}
It builds each module and writes
\texttt{.lake/build/leanquote/TranslationTiling/Abstract.leanquote}, inside the build directory.
From then on, \cs{leanquote} finds this export automatically and uses Lean's declaration ranges;
nothing changes in the document. The export also records the source as compiled. If the
\texttt{.lean} file is edited afterwards, the package warns and falls back to reading the
\texttt{.lean} file directly, so the document always shows the current code.
The exporter is a single Lean file (\texttt{lean/LeanQuoteExport.lean}) run with
\texttt{lake env lean --run}; nothing is added to the Lean project.
An export placed next to the \texttt{.lean} file is also found, and an export can be quoted
explicitly with \cs{leanfilequote}\verb|{Abstract.leanquote}{...}|.

\section{Engines and symbols}

The \texttt{leanquote} listings style typesets about a hundred of Lean's Unicode symbols
($\mathbb{N}\ \mathbb{Z}\ \mathbb{R}\ \to\ \forall\ \exists\ \sum\ \bigcup\ \le$ \dots, Greek letters, sub- and superscripts) as math, under
pdf\LaTeX{} through listings' \opt{literate} mechanism, and under Lua\LaTeX{} and Xe\LaTeX{} by
making these characters active inside listings. Other symbols can be added with \opt{literate}
in a style of your own (pdf\LaTeX), or by choosing a monospace font that contains them
(Lua\LaTeX, Xe\LaTeX).
The package loads \textsf{listings}, \textsf{xcolor}, \textsf{graphicx}, \textsf{amssymb},
and \textsf{tcolorbox} with its \texttt{skins} and \texttt{breakable}
libraries. It writes one temporary file \meta{jobname}\texttt{-leanquote-}\meta{n}\texttt{.lean}
per listing.

\section{Texlish}

\href{https://texlish.org}{Texlish} passes \LaTeX{} commands through, and its one-argument
syntax reads naturally:
\begin{latexsrc}
::use leanquote
\leanfile: Abstract.lean
\leanquote: main_theorem
\end{latexsrc}
With an optional argument, use braces: \verb|\leanquote[doc]{main_theorem}|.

\section{Limitations}

\begin{itemize}
\item Without an export, declaration boundaries follow the layout rules above; code that puts
  continuation lines in column~0 (other than the listed keywords) is cut short.
\item Namespaces are not tracked: two declarations with the same short name in different
  namespaces give an ambiguity warning; quote them with a name that distinguishes them.
\item The listings language is called \texttt{leanquote-lean}, so that a document's own
  \verb|\lstdefinelanguage{Lean4}| does not interfere with it.
\item The compiled check relies on file modification dates; copying a project without
  preserving them can cause false warnings or hide stale code.
\end{itemize}

\section{Related tools}

Several tools include code from source files in documents; as far as we know, none quotes Lean
declarations by name in an ordinary \LaTeX{} document or checks the quoted code against the
compiled module.
\begin{itemize}
\item \textbf{Marker-based inclusion in \LaTeX.}
  \textsf{catchfilebetweentags} (\verb|\ExecuteMetaData|), \textsf{listings} line ranges with
  marker comments, \textsf{piton} ranges, and \textsf{tagsnip} include the code between markers
  that the author adds to the source; \textsf{minted}'s \texttt{rangeregex} selects code by a
  regular expression, and Pygments has a Lean~4 lexer. These work for any language but need
  markers or patterns, and do not know Lean's declarations.
\item \textbf{Lean-native documents.} Verso and SubVerso, by the Lean FRO, include anchored
  regions of Lean modules in documents written in Verso, with Lean's own highlighting, and fail
  when a code block no longer matches its module; Verso's manual genre also produces \TeX.
  They are the natural choice for documents written in Lean; \textsf{leanquote} is for
  papers written in \LaTeX.
\item \textbf{Blueprints.} \textsf{leanblueprint}'s \verb|\lean{...}| links informal statements to
  Lean declarations by name, and \texttt{checkdecls} checks that the names exist;
  LeanArchitect generates blueprint text from annotated Lean declarations
  (\verb|\inputleannode{...}|). Both connect \LaTeX{} and Lean by name, without quoting code,
  and can be used together with \textsf{leanquote}.
\item \textbf{Other proof assistants.} Isabelle's document preparation and the
  \textsf{cartouches} tool (formerly \textsf{isasnips}), which names snippets after declarations;
  Agda's \LaTeX{} backend combined with \textsf{catchfilebetweentags}; and Alectryon for Rocq.
  There the checker itself generates the document, which keeps code and proofs in sync.
\item \textbf{Lean listings style.} Lean's documentation provides \texttt{lstlean.tex}, a
  \textsf{listings} definition of Lean with a Unicode table, which inspired ours.
\end{itemize}
The idea of including Lean snippets in papers to keep them in sync with the code was raised on
the Lean Zulip by Eric Wieser in 2022 (thread ``MathemaTeX'').

\section{Acknowledgments}

The package and this documentation were produced with the help of Claude Opus~5.5
(Anthropic), working under the authors' direction.

\section{License and history}

\textsf{leanquote} is released under the MIT License. It grew out of a position paper by the
authors on presenting the formal statements of machine-checked results to human readers.

\begin{description}
\item[0.8 (2026-10-25)] Trademark mark after the logo (\cs{leanquotelogomark}).
\item[0.7] \textsc{unverified} mark and \opt{verify} option; \cs{leanquotefile} and a safe \cs{leanfile}; listings language renamed \texttt{leanquote-lean}.
\item[0.6] Timestamps compared with \LaTeX3's \verb|\file_compare_timestamp|; related tools.
\item[0.5] Lua\LaTeX{} and Xe\LaTeX{} support; documentation.
\item[0.4] Renamed from \textsf{leandecl}; \cs{leanfile} and \cs{leanquote}.
\item[0.3] Optional exports used automatically when up to date.
\item[0.2] Direct reading of \texttt{.lean} files; compiled check.
\item[0.1] First version, based on exports.
\end{description}

\end{document}
